\subsubsection{Environmental Impact}

\textbf{Regarding PPP: Have you estimated the environmental impact of the project? }
The only resource directly affecting the environmental footprint is energy consumption. This resource is required throughout the entire project, as specified in Section 6.1.2. Both hardware resources and the shared workspace consume electricity continuously. However, sharing the workspace's lighting with other employees reduces its environmental impact.

\textbf{Regarding PPP: Did you plan to minimize its impact, for example, by reusing resources? }

The application uses Azure for software resources. Currently, the Azure resources required for application development cost less than traditional physical resources. This cost optimization results from significantly reduced resource consumption, as cloud resources eliminate the need for transportation and frequent physical replacements. Updates and fixes are performed virtually, further reducing hardware requirements and electronic waste.

\textbf{Regarding Useful Life: How is the problem that you want to address currently solved (state of the art)?}

The primary strategy for reducing resource usage is to develop the most efficient application possible, avoiding unnecessary resources and inefficient code that could overuse resources. This ensures that only essential resources are consumed each time employees use the application, limiting overall resource consumption throughout its lifecycle.

\textbf{Regarding Useful Life: How will your solution improve the environment with respect to other existing solutions? }

This application is tailored for the company's exclusive use and excludes any irrelevant functionality from the final product. It streamlines resource consumption by avoiding the unnecessary computation performed by other existing applications.

